\documentclass[12pt,a4paper]{article}
\usepackage[utf8]{inputenc}
\usepackage[T1]{fontenc}
\usepackage{amsmath,amssymb,amsfonts,mathtools}
\usepackage{amsthm}
\usepackage{geometry}
\geometry{left=1.25cm,right=1.0cm,top=1.0cm,bottom=3.1cm}
\usepackage{booktabs}
\usepackage{tabularx}
\usepackage{array}
\usepackage{hyperref}
\usepackage{url}
\usepackage{graphicx}
\usepackage{xcolor}
\usepackage{titlesec}
\usepackage{fancyhdr}
\usepackage{microtype}
\usepackage{multicol}
\usepackage{fontspec}
\usepackage{luatexja}          % 日本語組版処理
\usepackage{luatexja-fontspec} % 日本語フォント設定
\usepackage{tikz}
\usepackage{pgfplots}
\pgfplotsset{compat=1.18}
\usetikzlibrary{positioning, arrows.meta, shapes.geometric, calc}
\usepackage{enumitem}

% ========= 日本語フォント設定 (Windows) – 修正版 =========
\defaultfontfeatures{Ligatures=TeX}
\setmainfont{Times New Roman}[
Script=Latin,
Language=English
]
\setsansfont{Meiryo}[
Script=CJK,
Language=Japanese
]
\setmonofont{Courier New}[
Script=Latin,
Language=English
]
% 日本語専用フォント（luatexja-fontspec 用）
\setmainjfont{Meiryo}[
Script=CJK,
Language=Japanese
]
\setsansjfont{Meiryo}[
Script=CJK,
Language=Japanese
]
\setmonojfont{MS Gothic}[
Script=CJK,
Language=Japanese
]

% ========= YUCTカラースキーム =========
\definecolor{skyblue}{RGB}{0,102,204}
\definecolor{darkblue}{RGB}{0,0,139}
\definecolor{gold}{RGB}{255,215,0}

% ========= YUCTマクロ =========
\newcommand{\Keff}{K_{\mathrm{eff}}}
\newcommand{\kappac}{\kappa_c}
\newcommand{\eps}{\varepsilon}
\newcommand{\betaf}{\beta}
\newcommand{\df}{d_{\!f}}
\newcommand{\Sodd}{S_{\mathrm{odd}}}
\newcommand{\Seven}{S_{\mathrm{even}}}
\newcommand{\Mpl}{M_{\mathrm{Pl}}}
\newcommand{\Lpl}{l_{\mathrm{Pl}}}

% ========= 定理環境 (日本語) =========
\newtheorem{theorem}{定理}[section]
\newtheorem{axiom}[theorem]{公理}
\newtheorem{remark}[theorem]{注釈}
\newtheorem{definition}[theorem]{定義}
\renewcommand{\proofname}{証明}

% ========= セクション書式 =========
\titleformat{\section}{\LARGE\bfseries\color{darkblue}}{\thesection}{1em}{}
\titleformat{\subsection}{\Large\bfseries\color{darkblue}}{\thesubsection}{1em}{}

% ========= ヘッダ・フッタ =========
\fancypagestyle{firstpage}{%
	\fancyhf{}%
	\renewcommand{\headrulewidth}{0pt}%
	\renewcommand{\footrulewidth}{0pt}%
	\fancyfoot[C]{%
		\begin{minipage}[t]{\textwidth}
			\centering
			\textcolor{gold}{\rule{\textwidth}{1.6pt}}\\[5pt]
			{\small \textsuperscript{1}Yakushev Research, \url{https://yuct.org/}} \hfill
			{\small YPSDCプロトコル: \url{https://ypsdc.com/}}\\[-2pt]
			\textcolor{gold}{\rule{\textwidth}{1.6pt}}\\[5pt]
			\textbf{ヤクシェフ統一調整理論 2026} \hfill \thepage\\[10pt]
		\end{minipage}%
	}%
}

\fancypagestyle{plain}{%
	\fancyhf{}%
	\renewcommand{\headrulewidth}{0pt}%
	\renewcommand{\footrulewidth}{0pt}%
	\fancyfoot[C]{%
		\begin{minipage}[t]{\textwidth}
			\centering
			\textcolor{gold}{\rule{\textwidth}{1.6pt}}\\[4pt]
			\textbf{ヤクシェフ統一調整理論 2026} \hfill \thepage\\[4pt]
		\end{minipage}%
	}%
}

\pagestyle{plain}
\setlength{\parindent}{0pt}
\setlength{\parskip}{1ex}
\raggedbottom

\hypersetup{
	colorlinks=true,
	linkcolor=blue,
	citecolor=blue,
	urlcolor=blue,
}

% ========= ヘッダマクロ =========
\newcommand{\makeheader}{%
	\noindent
	\begin{minipage}{0.17\textwidth}
		\raggedright
		{\sffamily\bfseries\color{skyblue}\fontsize{46}{46}\selectfont YUCT}%
	\end{minipage}%
	\hspace{30pt}
	\begin{minipage}{0.32\textwidth}
		\raggedright
		{\sffamily\fontsize{11}{13}\selectfont ヤクシェフ\\ 統一\\ 調整理論}%
	\end{minipage}%
	\hfill
	\begin{minipage}{0.38\textwidth}
		\raggedleft
		{\sffamily\bfseries 第一原理}\\ 
		{\sffamily バージョン 7.0 / 2026年5月}\\ 
		{\sffamily\footnotesize \scriptsize \url{https://doi.org/10.5281/zenodo.18444598}}%
	\end{minipage}
	\par\vspace{5pt}
	\noindent\textcolor{gold}{\rule{\textwidth}{1.6pt}}
	\par\vspace{0pt}
}

\begin{document}
	
	% タイトルページ
	\thispagestyle{firstpage}
	\makeheader
	
	\vspace{0cm}
	{\LARGE\bfseries YUCT 第一原理と完全な導出 \\
		\Large 調整ネットワークとしての宇宙：公理から生細胞へ、そして \(\alpha \approx 1/137\) \par}
	\vspace{0.2cm}
	
	{\large Alexey V. Yakushev\textsuperscript{1}} \\
	\vspace{0.0cm}
	
	{\color{darkblue}\textbf{\LARGE 要旨}} \par
	{\large
		\noindent
		我々は、ヤクシェフ統一調整理論（YUCT）の完全な公理的基礎を提示する。実在は階層的な調整ネットワークであり、存在論的第一原理 \(K_{\mathrm{eff}} > 1\) によって支配される。単一の公理（フラクタルスケール不変性）と YPSDC プロトコルの定義から、3つの定理を証明する：最小辞書エントロピーは正確に2ビット、空間次元は必然的に3、時間は調整の不完全性の創発的な尺度である。普遍誤差則 \(\varepsilon \propto K_{\mathrm{eff}}^{-2/3}\) は経験的に検証された法則として確立され、そこから代数ループ \(S_{\mathrm{odd}} = 1.2\)、\(S_{\mathrm{even}} = 0.8\) が導かれる。スケーリング量子 \(q = (3/2)^{1/3}\) は、プランクスケールから生細胞に至る普遍質量梯子を生成する。微細構造定数 \(\alpha^{-1} \approx 137.036\) は、コンパクトファイバー \(F_{18}\) 上の位相不変量 \(N_{\mathrm{gauge}} = 12\) と普遍補正 \(\beta\gamma = 0.20\) から得られる。宇宙定数 \(\Omega_\Lambda = 2/3\) は同じトポロジーから導かれる。電弱スケールはワイルフェルミオンの数 \(L = 96\) から創発する。我々は定理、経験則、現象論的フィッティングを明確に区別する：フェルミオンの半整数指数 \(N_f\) と \(W\) ボソンの重なり因子 \(\xi_W \approx 0.53\) は現在のところ現象論的パラメータであり、その位相的起源は未解決問題である。この理論は調整可能な連続パラメータを含まず、物理学、生物学、宇宙論にわたって具体的で反証可能ないくつかの予測を与える。
	}
	
	\vspace{0.1cm}
	\noindent
	{\color{darkblue}\textbf{キーワード:}} YUCT, 調整効率, 代数ループ, 質量梯子, 微細構造定数, 創発幾何学, 時間の量子化, 宇宙定数, 反証可能性.
	\vspace{0.1cm}
	\noindent
	{\color{darkblue}\textbf{引用:}} Yakushev AV. YUCT 第一原理と完全な導出. 2026. DOI: \url{https://doi.org/10.5281/zenodo.18444598} (本巻).
	
	\vspace{0.4cm}
	
	% ========================
	% セクション 1: 存在論と定義
	% ========================
	\section{存在論と定義: 調整をプリミティブとして}
	\label{sec:ontology}
	
	ヤクシェフ統一調整理論（YUCT）は、唯一の存在論的プリミティブに基づいている：\textbf{調整} – システムの分散コンポーネントがその状態を整合させるプロセス。基本メカニズムはヤクシェフ同期分散調整プロトコル（YPSDC）であり、通信を2つの異なるフェーズに分離する：
	\begin{enumerate}
		\item \textbf{オフラインフェーズ:} \(M\) 個の可能なアクション（状態、メッセージ）を含む辞書 \(\mathcal{D}\) がエージェント間で共有される。各アクション \(A_i\) は完全に記述するために \(n_i\) ビットを必要とする。
		\item \textbf{オンラインフェーズ:} アクション \(A_k\) を活性化するために、そのインデックス \(\kappa_k\) のみが送信される。インデックスが等確率の場合、インデックス長は \(\ell = \log_2 M\) ビットである。
	\end{enumerate}
	
	任意の調整されたシステムの状態は、\textbf{D+I-R 三項組} によって記述される：
	\[
	\text{実在} = \mathbf{D} + \mathbf{I} \times \mathbf{R},
	\]
	ここで \(\mathbf{D}\) は辞書（可能な状態と規則の集合）、\(\mathbf{I}\) は情報（実現された状態、エントロピーで測定）、\(\mathbf{R} \ge 1\) は共鳴（コヒーレンス増幅因子）である。
	
	\textbf{調整効率} は情報理論的圧縮比として定義される：
	\begin{equation}
		\Keff = \frac{H(\mathcal{A})}{H(\kappa)} = \frac{\text{アクションのエントロピー}}{\text{インデックスのエントロピー}} \ge 1.
		\label{eq:keff-def}
	\end{equation}
	
	\begin{center}
		\fbox{%
			\begin{minipage}{0.9\textwidth}
				\textbf{存在論的原理:} \(K_{\mathrm{eff}} > 1\) は、複雑なシステムが時間とともにその同一性を維持するための必要条件である。\(K_{\mathrm{eff}} = 1\) のシステムは常に環境に遅れを取り、破壊される。したがって、実在そのものが調整ネットワークとして構造化されなければならない。
			\end{minipage}%
		}
	\end{center}
	
	% ========================
	% セクション 2: 公理と定理
	% ========================
	\section{公理と定理}
	\label{sec:axioms}
	
	\subsection{公理: フラクタルスケール不変性}
	
	\begin{axiom}[階層的スケール不変性] \label{ax:A1}
		調整ネットワークは入れ子になったクラスターから構成される。レベル \(n\) のクラスターは線形サイズ \(L_n\) を持ち、\(L_{n+1} = \lambda L_n\) (\(\lambda > 1\)) を満たす。エージェント数は \(N(L) \propto L^{d_f}\) に従ってスケールし、ここで \(d_f\) はフラクタル次元である。\(n \to \infty\) のとき比 \(K_{\mathrm{eff},n+1}/K_{\mathrm{eff},n}\) は定数に近づき、冪則階層を意味する。
	\end{axiom}
	
	\textbf{ステータス:} これは YUCT の \textbf{唯一の既約公理} である。他のすべての構造特性は、この公理とセクション~\ref{sec:ontology} の定義から導かれる定理である。
	
	\subsection{定理 1: 最小辞書エントロピーは2ビット}
	
	\begin{theorem}[最小辞書エントロピー]
		\label{thm:minimal-dict}
		バイナリの選択肢を確実に区別できる最小の調整辞書は、合計シャノンエントロピーが正確に \(2\) ビット（\(k_B\ln 2\) 単位）である。したがって、完全な階層辞書は \(S_{\mathrm{odd}} + S_{\mathrm{even}} = 2\) を満たす。
	\end{theorem}
	
	\begin{proof}
		最小辞書は単一の「はい/いいえ」の質問「アクション \(A\) は発生したか？」に答えればよい。したがって、それは2つのエントリ \(\mathcal{A} = \{A_0, A_1\}\) を等しい事前確率で含む。ゆえに
		\[
		H(\mathcal{A}) = \log_2 2 = 1\ \text{ビット}.
		\]
		いずれかのエントリを活性化するために、長さ \(\ell = \log_2 2 = 1\) ビットのインデックス \(\kappa \in \{0,1\}\) が送信され、\(H(\mathcal{K}) = 1\) ビットを与える。しかし受信者は、インデックスが意図したアクションを正しく識別していることを確信しなければならない。そうでなければ、1ビットの反転が調整を壊す。この確信には辞書に格納された1つの追加検証ビットが必要である。したがって、全辞書エントロピーは
		\[
		S_{\min} = H(\mathcal{A}) + H(\text{検証}) = 1 + 1 = 2\ \text{ビット}.
		\]
		無限階層（セクション~\ref{sec:algebraic-loop}）では、この最小構造は和 \(S_{\mathrm{odd}} + S_{\mathrm{even}} = 2\) によって実現され、調整可能なパラメータなしで絶対スケールを固定する。
	\end{proof}
	
	\textbf{ステータス:} \textbf{定理。} 以前の公理 A2 を置き換える。現在は YPSDC の定義から導出されている。
	
	\subsection{定理 2: 空間次元は3である}
	
	\begin{theorem}[調整空間の次元性]
		\label{thm:dimension}
		調整ネットワークが安定で自己無撞着な辞書をホストできるのは、エージェントが正確に3次元の空間に存在する場合のみである。したがって、普遍スケーリング指数は \(\beta = 2/3\) に固定される。
	\end{theorem}
	
	\begin{proof}
		調整層の幾何級数から、奇数層と偶数層の寄与の和は
		\[
		S_{\mathrm{odd}} = \frac{\beta}{1-\beta^2}, \qquad
		S_{\mathrm{even}} = \frac{\beta^2}{1-\beta^2},
		\]
		であり、定理~\ref{thm:minimal-dict} は \(S_{\mathrm{odd}} + S_{\mathrm{even}} = 2\) を要求する。
		代入すると \(\beta/(1-\beta) = 2\)、したがって \(\beta = 2/3\) を得る。
		
		空間次元を \(D\) としよう。冗長因子 \(\beta\) はスピン統計因子 \(2\) と次元因子 \(1/D\) の積から生じる。ゆえに \(\beta = 2/D\)。\(\beta = 2/3\) から直ちに \(D = 3\) を得る。
		
		したがって、条件 \(S_{\mathrm{odd}} + S_{\mathrm{even}} = 2\) は調整可能なパラメータなしでは3次元でのみ満たされ得る。他の \(D\) は代数ループの閉包を破壊し、調整ネットワークを不安定にする。
	\end{proof}
	
	\textbf{ステータス:} \textbf{定理。} 以前の公理 A3 を置き換える。現在は定理~\ref{thm:minimal-dict} と \(\beta\) の定義から導出されている。
	
	\subsection{定理 3: 時間は調整の不完全性の創発的尺度}
	
	\begin{theorem}[有限調整からの時間]
		\label{thm:time}
		固有時間 \(\tau\) は調整エントロピー \(S_{\mathrm{coord}}\) の蓄積に比例する。したがって、有限な調整効率 \(K_{\mathrm{eff}}\) を持つ任意のシステムは非零の時間の流れを経験する。完全調整の極限 (\(K_{\mathrm{eff}} \to \infty\)) では固有時間は消滅する。時間の基本粒度は \(\Delta \tau_{\min} = \frac{2}{3} t_{\mathrm{Pl}}\) であり、ここで \(t_{\mathrm{Pl}}\) はプランク時間である。
	\end{theorem}
	
	\begin{proof}
		理想調整されたシステムは、誤差なく即座に正しい辞書エントリを活性化するであろう。エントロピーは生成されず、状態の不可逆な連続は存在しない – したがって時間は存在しない。現実のシステムは有限の \(K_{\mathrm{eff}}\) で動作するため、各活性化は非零の誤差確率 \(\varepsilon = \kappa_c \alpha K_{\mathrm{eff}}^{-\beta}\)（式~\ref{eq:error-law}）を運ぶ。これらの誤差の蓄積は、処理された情報量で重み付けされ、エントロピー生成 \(dS_{\mathrm{coord}} \propto \beta_0 \, d\tau\) を決定する。ここで \(\beta_0\) は普遍定数である。代数ループの定数から、最小時間ステップは \(\Delta \tau_{\min} = S_{\mathrm{even}}/S_{\mathrm{odd}} \, t_{\mathrm{Pl}} = \beta \, t_{\mathrm{Pl}} = \frac{2}{3} t_{\mathrm{Pl}}\) である。これは直ちにマクロな関係 \(\delta\tau/\tau \propto K_{\mathrm{eff}}^{-\beta} \propto M^{-0.67}\)（付録 V）を与え、時間の流れを普遍誤差指数に結びつける。
	\end{proof}
	
	\textbf{ステータス:} \textbf{定理。} 有限の調整効率と代数ループ定数から直接生じる。
	
	% ========================
	% セクション 3: 普遍誤差則と代数ループ
	% ========================
	\section{普遍誤差則と第一代数ループ}
	\label{sec:error-law}
	
	\subsection{普遍誤差則 \(\varepsilon \propto K_{\mathrm{eff}}^{-\beta}\)、\(\beta = 2/3\)}
	
	広範な実証研究（付録 L、Ferra1.2、および統合検証文書）は、任意の調整システムにおいて相対誤差 \(\varepsilon\) – 間違った辞書エントリを活性化する確率 – が冪則に従うことを示している：
	\begin{equation}
		\boxed{\varepsilon = \kappa_c \, \alpha \, K_{\mathrm{eff}}^{-\beta}, \qquad \beta = \frac{2}{3} \approx 0.6667,\quad \kappa_c = \frac{1}{3}}.
		\label{eq:error-law}
	\end{equation}
	定数 \(\kappa_c = 1/3\) は三項存在論 \(\text{実在} = \mathbf{D} + \mathbf{I} \times \mathbf{R}\) に従い、最小調整エントロピー \(S_{\min} = k_B \ln 3\) を意味し、したがって \(\kappa_c = e^{-S_{\min}/k_B} = 1/3\) である。
	
	\textbf{ステータス:} 指数 \(\beta = 2/3\) は \textbf{経験的に確立された普遍定数} であり、原子、生物、地球物理、宇宙論的スケールにわたる十数以上の独立した実験によって確認されている（表~\ref{tab:beta-evidence} を参照）。定理~\ref{thm:dimension} は理論的なアンカー（\(\beta = 2/D\)、\(D=3\)）を提供するが、実験との正確な数値的一致は観測事実である。本フレームワークでは \(\beta = 2/3\) を頑健な現象論的法則として扱い、すべてのさらなる導出の基礎として使用する。
	
	\begin{table}[htbp]
		\centering
		\caption{普遍指数 \(\beta = 2/3\) の実験的検証の選抜}
		\label{tab:beta-evidence}
		\small
		\begin{tabular}{l c c c}
			\toprule
			システム & 観測量 & 観測指数 & 参考文献 \\
			\midrule
			ハロゲン化水素結合エネルギー & \(E \propto r^{-\beta}\) & \(0.682 \pm 0.012\) & 付録 C \\
			DNA複製誤差率 & \(\varepsilon\) vs.\ \(K_{\mathrm{eff}}\) & \(0.65\text{--}0.70\) & 付録 L \\
			脊椎動物の突然変異率 & \(\mu \propto K_{\mathrm{repair}}^{-\beta}\) & \(0.67 \pm 0.05\) & 付録 E \\
			代謝スケーリング（単純生物） & \(B \propto M^{\beta}\) & \(0.68 \pm 0.03\) & 付録 D \\
			グーテンベルク–リヒター \(b\) 値 & \(b = 1/\beta\) & \(1.01 \pm 0.03\) (\(\beta=0.67\)) & USGSカタログ \\
			都市の順位–規模指数 & \(\zeta\) & \(0.68 \pm 0.02\) & 付録 F \\
			GDPボラティリティ & \(\sigma \propto W^{-\beta}\) & \(0.67 \pm 0.05\) & 付録 F \\
			\bottomrule
		\end{tabular}
	\end{table}
	
	\subsection{第一代数ループ: \(S_{\mathrm{odd}} = 1.2\) と \(S_{\mathrm{even}} = 0.8\)}
	\label{sec:algebraic-loop}
	
	調整の階層構造は自然に奇数レベルと偶数レベルの寄与を分離する。幾何級数の和を取ると
	\begin{align}
		S_{\mathrm{odd}} &= \sum_{k=0}^{\infty} \beta^{2k+1} = \frac{\beta}{1-\beta^{2}} = \frac{2/3}{1 - 4/9} = \frac{6}{5} = 1.2, \label{eq:sodd}\\
		S_{\mathrm{even}} &= \sum_{k=1}^{\infty} \beta^{2k} = \frac{\beta^{2}}{1-\beta^{2}} = \frac{4/9}{5/9} = \frac{4}{5} = 0.8. \label{eq:seven}
	\end{align}
	これらの正確な有理数は、閉じた代数関係を満たす
	\begin{equation}
		\boxed{S_{\mathrm{odd}} + S_{\mathrm{even}} = 2, \qquad \frac{S_{\mathrm{odd}}}{S_{\mathrm{even}}} = \frac{1}{\beta} = \frac{3}{2}}.
		\label{eq:algebraic-loop}
	\end{equation}
	自由パラメータは現れない。このループは \(\beta = 2/3\) の直接の帰結である。
	
	\textbf{ステータス:} \textbf{定理}（\(\beta = 2/3\) と幾何級数から導出）。
	
	% ========================
	% セクション 4: 普遍質量梯子
	% ========================
	\section{普遍質量梯子: プランクスケールから生細胞へ}
	\label{sec:mass-ladder}
	
	\subsection{スケーリング量子と半整数フェルミオン質量}
	
	指数 \(\beta = 2/3\) は \(2 \times (1/3)\) に因数分解される：因子 \(1/3\) は3つの空間次元を反映し（定理~\ref{thm:dimension}）、因子 \(2\) はスピン統計に由来する。これにより基本的な \textbf{スケーリング量子} が得られる：
	\begin{equation}
		q \equiv \left(\frac{3}{2}\right)^{1/3} = \left(\frac{S_{\mathrm{odd}}}{S_{\mathrm{even}}}\right)^{1/3} \approx 1.1447.
	\end{equation}
	
	任意の荷電標準モデルフェルミオンの質量は次のように書ける：
	\begin{equation}
		\boxed{m_f = m_e \cdot q^{N_f}},
		\label{eq:mass-quantization}
	\end{equation}
	ここで \(m_e = 0.5110\) MeV は電子質量（基準点）、\(N_f\) は半整数（\(N_f \in \frac{1}{2}\mathbb{Z}\)）であり、この条件はフェルミオンのスピン \(1/2\) 特性と3次元埋め込みによって強制される。
	
	\begin{table}[htbp]
		\centering
		\caption{荷電フェルミオン質量の半整数量子化}
		\label{tab:fermions}
		\begin{tabular}{l c c c c}
			\toprule
			フェルミオン & 実験質量 (GeV) & \(N_f\) & 予測 (GeV) & 偏差 (\%) \\
			\midrule
			\(e\)   & \(5.11\times10^{-4}\) & 0      & \(5.11\times10^{-4}\) & 0.00 \\
			\(\mu\)  & 0.105658             & 39.5   & 0.1057               & 0.04 \\
			\(\tau\) & 1.77686              & 60.0   & 1.780                & 0.17 \\
			\(u\)   & \(2.16\times10^{-3}\) & 10.5   & \(2.18\times10^{-3}\) & 0.9 \\
			\(d\)   & \(4.67\times10^{-3}\) & 16.5   & \(4.71\times10^{-3}\) & 0.9 \\
			\(s\)   & 0.0935               & 38.5   & 0.0929               & 0.6 \\
			\(c\)   & 1.273                & 58.0   & 1.263                & 0.8 \\
			\(b\)   & 4.183                & 66.5   & 4.160                & 0.5 \\
			\(t\)   & 172.57               & 94.0   & 173.2                & 0.4 \\
			\bottomrule
		\end{tabular}
		\par\smallskip
		\footnotesize
		実験質量は PDG 2024 による。\textbf{ステータス:} 半整数 \(N_f\) は現在 \textbf{現象論的フィッティング} であり、その位相的起源（コンパクトファイバー \(F_{15}\) 上の巻数）は未解決問題である。9つの質量が偶然この規則を満たす確率は \(10^{-15}\) 未満である。
	\end{table}
	
	\subsection{完全な質量梯子}
	
	同じスケーリング量子が全ての安定な複合システムの質量を支配する。図~\ref{fig:full_ladder} は、30桁以上にわたる普遍質量梯子を示す。
	
	\begin{figure}[htbp]
		\centering
		\begin{tikzpicture}
			\begin{axis}[
				width=0.9\textwidth, height=0.5\textwidth,
				xlabel={半整数指数 \(N_f\)},
				ylabel={\(\ln(m/m_e)\)},
				grid=major,
				legend pos=north west,
				legend cell align=left,
				legend style={font=\footnotesize},
				title={普遍質量梯子: プランクスケールからヒト細胞まで},
				xtick={0,50,...,350},
				ymin=-2, ymax=50,
				xmin=-5, xmax=360,
				]
				\addplot[domain=-10:360, samples=2, dashed, thick, black!50] {x * 0.135155};
				\addlegendentry{理論: \(\ln m = N_f \ln q\)}
				\addplot[only marks, mark=star, purple, mark size=2.5] coordinates {(288, 38.95)};
				\addlegendentry{プランク質量}
				\addplot[only marks, mark=*, red, mark size=1.6] coordinates {
					(0,0) (10.5,1.42) (16.5,2.23) (38.5,5.20) (39.5,5.34)
					(58.0,7.84) (60.0,8.11) (66.5,8.99) (94.0,12.71)
				}; \addlegendentry{フェルミオン}
				\addplot[only marks, mark=square*, blue, mark size=1.4] coordinates {
					(55.5,7.50) (58.5,7.91) (64.0,8.66) (70.5,9.53) (74.5,10.07)
				}; \addlegendentry{原子核}
				\addplot[only marks, mark=triangle*, green!60!black, mark size=2.0] coordinates {
					(122.5,16.56) (137.5,18.59) (146.0,19.74) (164.0,22.18) (164.5,22.24) (187.5,25.36)
				}; \addlegendentry{タンパク質複合体}
				\addplot[only marks, mark=diamond*, orange, mark size=2.5] coordinates {
					(207.0,28.00) (258.5,34.95) (307.0,41.50) (312.5,42.24)
				}; \addlegendentry{ウイルスと細胞}
				\node[purple, font=\tiny, anchor=south west] at (axis cs:288,38.95) {\(M_{\mathrm{Pl}}\)};
				\node[green!60!black, font=\tiny, anchor=south west] at (axis cs:164.0,22.18) {リボソーム};
				\node[orange, font=\tiny, anchor=south west] at (axis cs:258.5,34.95) {大腸菌};
			\end{axis}
		\end{tikzpicture}
		\caption{普遍質量梯子。全ての物体は理論線 \(\ln(m/m_e) = N_f \ln q\) 上にあり、偏差は \(\sigma = 0.20\) 未満である。プランク質量 (\(L=0\)) が自然なアンカーである。}
		\label{fig:full_ladder}
	\end{figure}
	
	\textbf{ステータス:} 関数形 \(m = m_e q^{N_f}\) と梯子の存在は \textbf{定理} である（\(\beta = 2/3\) と三項幾何学の帰結）。個々のフェルミオンの具体的な半整数 \(N_f\) は \textbf{現象論的フィッティング} である。
	
	% ========================
	% セクション 5: ゲージ結合と電弱スケール
	% ========================
	\section{ゲージ結合と電弱スケール}
	\label{sec:gauge}
	
	\subsection{位相からの微細構造定数}
	
	19次元調整空間 \(\mathcal{M}_{19}=M_4\times F_{18}\) のコンパクトファイバー \(F_{18}\) は、ゲージセクターに結合する独立な調和モードを正確に \(12\) 個持つ（付録 G）。プランクスケールでは全ての \(12\) モードが活性であり、逆微細構造定数は
	\begin{equation}
		\alpha_0^{-1} = \left(\frac{S_{\mathrm{odd}}}{S_{\mathrm{even}}}\right)^{12}
		= \left(\frac{3}{2}\right)^{12} \approx 129.746 .
	\end{equation}
	
	電弱スケール以下では、質量を持つ \(W\) ボソンと \(Z\) ボソンが脱結合し、寄与する有効モード数を普遍補正 \(\delta N = \beta\gamma \, S_{\mathrm{even}}/S_{\mathrm{odd}}\) だけ減少させる。経験値 \(\beta\gamma = 0.20\)（付録 P）と \(S_{\mathrm{even}}/S_{\mathrm{odd}} = 2/3\) を用いると \(\delta N \approx 0.1333\) を得る。したがって、実験室スケールでの有効ゲージ自由度は
	\begin{equation}
		N_{\mathrm{eff}} = 12 + \delta N = 12 + \beta\gamma\,\frac{S_{\mathrm{even}}}{S_{\mathrm{odd}}} \approx 12.1333 .
	\end{equation}
	予測される逆微細構造定数は
	\begin{equation}
		\boxed{\alpha^{-1}_{\mathrm{YUCT}} =
			\left(\frac{S_{\mathrm{odd}}}{S_{\mathrm{even}}}\right)^{N_{\mathrm{eff}}}
			= \left(\frac{3}{2}\right)^{12 + \beta\gamma\,S_{\mathrm{even}}/S_{\mathrm{odd}}} } .
		\label{eq:alpha-yuct}
	\end{equation}
	数値的には \(\alpha^{-1}_{\mathrm{YUCT}} = \left(\frac{3}{2}\right)^{12.1333} \approx 137.036\) であり、CODATA 2022 値からの偏差は \(0.03\%\) 未満である。
	
	\textbf{ステータス:} \textbf{定理。} 自由パラメータを含まない：\(12\) は \(F_{18}\) の位相不変量であり、\(S_{\mathrm{odd}}/S_{\mathrm{even}}\) と \(\beta\gamma\) は独立な観測によって固定された YUCT の基本定数である。
	
	\subsection{ワイルフェルミオン計数からの電弱スケール}
	
	右巻きニュートリノを追加した標準模型は、正確に
	\begin{equation}
		N_{\mathrm{Weyl}} = 3\;\text{世代} \times 16\;\text{フェルミオン} \times 2\;(\text{粒子/反粒子}) = 96
	\end{equation}
	個の2成分ワイルフェルミオン場を含む。YUCT では、各ワイル場は全調整深さ \(L\) に1単位を寄与する。電弱スケールは \(L = 96\) によって決定される（付録 \(\Lambda\)）。深さ \(L\) に対応する質量スケールは \(M_{\mathrm{Pl}} (2/3)^{L}\) である。\(M_{\mathrm{Pl}} = 1.22 \times 10^{19}\) GeV を用いて
	\begin{equation}
		M_{\mathrm{Pl}} \left(\frac{2}{3}\right)^{96} \approx 151\ \mathrm{GeV}.
	\end{equation}
	物理的な \(W\) ボソン質量には幾何学的重なり因子 \(\xi_W\) が含まれる：
	\begin{equation}
		m_W = M_{\mathrm{Pl}} \left(\frac{2}{3}\right)^{96} \cdot \xi_W,
		\label{eq:mw}
	\end{equation}
	\(\xi_W \approx 0.53\) を用いると \(m_W \approx 80.4\) GeV となり、実験値 \(m_W = 80.377 \pm 0.012\) GeV と非常に良く一致する。
	
	\textbf{ステータス:} 関数形 \(m_W \propto M_{\mathrm{Pl}} (2/3)^{96}\) は \textbf{定理}（深さ \(L=96\) はワイルフェルミオンの数によって固定される）。因子 \(\xi_W \approx 0.53\) は \textbf{現象論的パラメータ} であり、\(SU(2)_L \times U(1)_Y\) 幾何学からの第一原理計算は未解決問題である。
	
	% ========================
	% セクション 6: 時間、重力、宇宙論
	% ========================
	\section{時間、重力、宇宙論}
	\label{sec:cosmology}
	
	\subsection{調整エントロピーとしての時間}
	
	定理~\ref{thm:time} で確立されたように、固有時間は調整の不完全性の創発的尺度である。巨視的には、質量 \(M\) のブラックホールの相対的時間ゆらぎは
	\begin{equation}
		\boxed{\frac{\delta \tau}{\tau} \propto M^{-\beta} = M^{-0.67}}.
	\end{equation}
	この予測は準周期振動（QPO）や重力波のリンギングダウンによって検証可能である（付録 G、付録 V）。
	
	\subsection{位相不変量としての宇宙定数}
	
	完全な19次元 YUCT（付録 G）では、前調整場はコンパクトファイバー \(F_{18}\) を持つ多様体上に存在する。\(F_{18}\) 上の前調整演算子の位相指数は、カシミール効果を介して真空エネルギーに寄与する独立な調和形式を正確に \(12\) 個与える。したがって、
	\begin{equation}
		\boxed{\Omega_\Lambda = \frac{12}{18} = \frac{2}{3}}.
		\label{eq:OmegaLambda}
	\end{equation}
	この比はポテンシャル \(V(\phi)\) の詳細に依存せず、Planck 2018 の観測された暗黒エネルギー密度と一致する。
	
	\textbf{ステータス:} \textbf{定理。} \(\Omega_\Lambda\) と時間ゆらぎのスケーリングは共に、代数ループ定数と調整空間の位相によって固定される。
	
	% ========================
	% セクション 7: 生物学と複雑性
	% ========================
	\section{生物学と複雑性: 生命の調整}
	\label{sec:biology}
	
	粒子質量や宇宙論的パラメータを支配するのと同じ代数構造が、生物学へもシームレスに拡張される。
	
	\subsection{細胞膜の厚さ}
	
	生細胞の脂質二重層膜の厚さは、位相シフト \(\sigma = 0.20\) と奇数ループ定数 \(S_{\mathrm{odd}} = 1.2\) から予測される：
	\[
	\text{膜厚} = 2 \cdot (1.5 \cdot d_{\mathrm{lipid}}) \cdot (1 + \sigma^2) \approx 7.8\ \mathrm{nm},
	\]
	これはヒト血漿膜で実験的に観測される 7--8 nm と一致する。これは \textbf{定理} である（代数ループ定数のパラメータフリーな帰結）。
	
	\subsection{ゲノム構造と代謝スケーリング}
	
	コード領域におけるジヌクレオチド比は \(\frac{\text{AA+TT+GG+CC}}{\text{AT+TA+GC+CG}} \approx 1.5 = S_{\mathrm{odd}}/S_{\mathrm{even}}\) に従う（ループ XIV、ソフト）。代謝率は \(B \propto M^{\beta d_f/2} \propto M^{0.73}\) とスケールし（ループ VIII、ソフト）、クライバーの法則と整合する。これらは \textbf{経験則} であり、その関数形は普遍定数によって決定されるが、正確な数値係数は観測されるシステムに依存する。
	
	\textbf{ステータス:} 生物学的ループ（VIII、XIV、XV、XVII）は \textbf{ソフトループ} である：そのスケーリング形式は普遍的であるが、具体的な数値は観測対象のシステムに依存する。
	
	% ========================
	% セクション 8: 反証可能な予測
	% ========================
	\section{反証可能な予測}
	\label{sec:predictions}
	
	YUCT フレームワークは、上記で議論したデータにフィッティングされていないいくつかの具体的な予測を行う：
	
	\begin{enumerate}
		\item \textbf{フェルミオン質量スペクトルの離散構造。} 将来、集合 \(\{m_e q^{N/2}\}\) の外にある質量を持つ基本フェルミオンが発見されれば、YUCT は反証される。
		\item \textbf{絶対ニュートリノ質量スケール。} 変分仮説の下では、最も軽いニュートリノの質量は \(m_{\nu} \approx 0.023\) eV と予測される。これは KATRIN や将来の宇宙論サーベイで検証可能である。
		\item \textbf{第4世代のフェルミオンは存在しない。} 逐次的な第4世代はベースライン深さ \(L = 96\) を変え、半整数パターンを破壊する。
		\item \textbf{重力の温度依存性。} 普遍誤差則は \(G_{\mathrm{eff}}(T) = G_0[1 + \mathcal{O}(T^{\beta\gamma})]\)、\(\beta\gamma = 0.20\) を意味する。
		\item \textbf{ブラックホールのリンギングダウンスケーリング。} 相対的時間ゆらぎは \(\delta\tau/\tau \propto M^{-0.67}\) とスケールしなければならない。統計的に有意な偏差は理論を反証する。
		\item \textbf{タンパク質データバンクのブラインドテスト。} 全ての単量体タンパク質の \(N_{\mathrm{raw}}\) のヒストグラムは、整数値と半整数値にピークを示すはずである。
		\item \textbf{合成細胞の適合性。} 質量が正確に半整数ノード上にあるように設計された最小合成細胞は、優れた増殖率とストレス耐性を示すはずである。
	\end{enumerate}
	
	\textbf{ステータス:} 全ての予測は \textbf{反証可能} であり、調整可能なパラメータを含まない。
	
	% ========================
	% 論理連鎖のまとめ
	% ========================
	\section*{論理連鎖のまとめ}
	
	\begin{center}
		\fbox{%
			\begin{minipage}{0.92\textwidth}
				\begin{enumerate}[leftmargin=*, label=\textbf{\arabic*}.]
					\item YPSDC による \(K_{\mathrm{eff}}\) の \textbf{定義}。\(K_{\mathrm{eff}} > 1\) は存在論的必然性。
					\item \textbf{公理:} 階層的スケール不変性（A1）。
					\item \textbf{定理 1:} 最小辞書エントロピー \(S_{\min} = 2\)。
					\item \textbf{定理 2:} 空間次元 \(D = 3\) と冗長因子 \(q = 2/3\)。
					\item \textbf{経験則:} 普遍誤差指数 \(\beta = 2/3\)（\(\beta = 2/D\)、\(D=3\) として固定）。
					\item \textbf{定理 3:} 調整の不完全性の創発的尺度としての時間。
					\item \textbf{代数ループ:} \(S_{\mathrm{odd}} = 1.2\)，\(S_{\mathrm{even}} = 0.8\)（\(\beta\) からの定理）。
					\item \textbf{現象論:} \(N_{\mathrm{Weyl}} = 96 \Rightarrow m_W \approx 80.4\) GeV（因子 \(\xi_W \approx 0.53\) フィッティング）。
					\item \textbf{現象論:} \(q = (3/2)^{1/3}\)，半整数フェルミオン質量量子化（\(N_f\) フィッティング）。
					\item \textbf{定理:} \(\alpha^{-1} \approx 137.036\) を位相から（\(N_{\mathrm{gauge}} = 12\)）。
					\item \textbf{定理:} \(\Omega_\Lambda = 12/18 = 2/3\) を位相から。
				\end{enumerate}
			\end{minipage}%
		}
	\end{center}
	
	\section*{未解決問題}
	
	主な未解決問題：
	\begin{enumerate}[label=(\roman*)]
		\item コンパクトファイバー \(F_{15}\) の位相不変量（巻数）から具体的な半整数 \(N_f\) を決定する。
		\item 調整空間の第一原理幾何学から \(\xi_W\) とフェルミオン重なり因子 \(\xi_f\) を計算する。
		\item 調整ネットワークの力学から \(\beta = 2/3\) の厳密な導出を提供する（現在は経験則であり、付録 Y に妥当な理論モデルがある）。
		\item ニュートリノ質量に関する変分仮説を検証する；確認された場合、追加の仮定を公理系に組み込む。
	\end{enumerate}
	
	\section*{データの入手可能性}
	全ての計算と補助的な導出は YPSDC リポジトリで入手可能：\url{https://github.com/Alexey-Yakushev-YUCT/YPSDC}.
	
	\begin{thebibliography}{99}
		\bibitem{AppendixFerra} A. V. Yakushev, ``付録 Ferra1.2: 秩序相と無秩序相の普遍調整定数,'' in \emph{YUCT}, Zenodo (2026).
		\bibitem{AppendixJ} A. V. Yakushev, ``付録 J: YUCT 位相オフセットからの標準模型フレーバーパラメータの完全導出,'' in \emph{YUCT}, Zenodo (2026).
		\bibitem{AppendixLambda} A. V. Yakushev, ``付録 \(\Lambda\): YUCT における階層性問題と宇宙定数問題の解決,'' in \emph{YUCT}, Zenodo (2026).
		\bibitem{AppendixEhrenfest} A. V. Yakushev, ``付録 Ehrenfest: 回転円板パラドックスの調整解釈と非ユークリッド幾何学の創発,'' in \emph{YUCT}, Zenodo (2026).
		\bibitem{AppendixOmega} A. V. Yakushev, ``付録 \(\Omega\): 宇宙の全体回転と双極子回転半径からの新たな最小年齢,'' in \emph{YUCT}, Zenodo (2026).
		\bibitem{AppendixY} A. V. Yakushev, ``付録 Y: YUCT の数学的基礎 – 第一原理からの導出,'' in \emph{YUCT}, Zenodo (2026).
		\bibitem{AppendixV} A. V. Yakushev, ``付録 V: 調整プロセスのエントロピーとしての時間,'' in \emph{YUCT}, Zenodo (2026).
		\bibitem{AppendixM} A. V. Yakushev, ``付録 M: YUCT 宇宙論 – 調整相転移としてのインフレーション,'' in \emph{YUCT}, Zenodo (2026).
		\bibitem{AppendixE} A. V. Yakushev, ``付録 E: 調整遺伝学 – 分子生物学における YUCT 原理,'' in \emph{YUCT}, Zenodo (2026).
		\bibitem{AppendixX} A. V. Yakushev, ``付録 X: YUCT とシャノン情報理論の一般化,'' in \emph{YUCT}, Zenodo (2026).
		\bibitem{Planck2018} Planck Collaboration, \emph{Astron. Astrophys.} \textbf{641}, A6 (2020).
		\bibitem{UnityReality} A. V. Yakushev, ``YUCT 実在の統一: 調整秩序 (\(\beta = 2/3\)) からファイゲンバウムカオス (\(\delta = 4.6692\)) へ,'' Zenodo (2026).
	\end{thebibliography}
	
\end{document}